{{دیوان
| عنوان = تصوّف
| ماخذ = ویکی ماخذ
| ماخذ_لنک = https://ur.wikisource.org/wiki/%D8%AA%D8%B5%D9%88%D9%91%D9%81
| مقابلہ = اقبال سائبر لائبریری (اقبال اکادمی پاکستان)
| مقابلہ_لنک = https://www.iqbalcyberlibrary.net
}}
<poem>
یہ حکمتِ ملکوتی، یہ علمِ لاہُوتی
حرم کے درد کا درماں نہیں تو کچھ بھی نہیں

یہ ذکرِ نیم شبی، یہ مراقبے، یہ سُرور
تری خودی کے نگہباں نہیں تو کچھ بھی نہیں

یہ عقل، جو مہ و پرویں کا کھیلتی ہے شکار
شریکِ شورشِ پنہاں نہیں تو کچھ بھی نہیں

خِرد نے کہہ بھی دیا ’لااِلہ‘ تو کیا حاصل
دل و نگاہ مسلماں نہیں تو کچھ بھی نہیں

عجب نہیں کہ پریشاں ہے گفتگو میری
فروغِ صبح پریشاں نہیں تو کچھ بھی نہیں
</poem>

٭ ریاض منزل (دولت کدہ سرراس مسعود) بھوپال میں لکھے گئے
